\section{Augmented States}
\label{section_augmented_states}

Our hypothesis is that the control input will remain constant for $N_d$ consecutive cycles. To 
incorporate this into our system model, we will expand the prediction equation over these $N_d$ 
steps and then redefine the terms for simplification.

Starting with the standard discrete-time linear prediction equation \refeq{cap3-dynamic-equation}:

\noindent
\begin{equation}
    x[k + 1] = A x[k] + B u[k] \nonumber
\end{equation}

We can iteratively expand this equation for several steps, assuming the control 
input $u$ remains constant at $u\big[k\big]$ for the next $N_d - 1$ steps:

\noindent
\begin{align}
    x\big[k+1\big] &= A x\big[k\big] + B u\big[k\big] \nonumber \\
    x\big[k+2\big] &= A^2 x\big[k\big] + AB u\big[k\big] + B u\big[k\big] \nonumber \\
    x\big[k+3\big] &= A^3 x\big[k\big] + A^2B u\big[k\big]+ AB u\big[k\big] + B u\big[k\big] \nonumber \\
    &\vdots \nonumber \\
    x\big[k+N_d\big] &= A^{N_d} x\big[k\big] + \big(A^{N_{d-1}} + A^{N_{d-2}} + \dots +  A^0 \big)Bu\big[k\big]
\end{align}

To simplify the notation, let's define new matrices $A_b$ and $B_b$ as:
\noindent
\begin{align}
    A_b &= A^{N_d} \\
    B_b &= \big(A^{N_{d-1}} + A^{N_{d-2}} + \dots +  A^0 \big)B
\end{align}

Furthermore, let's shift the time index by blocks of $N_d$ cycles. This implies the following relationship between 
the original time index $k$ and the new time index $k'$:

\noindent
\begin{align*}
    k' - 1 &= k - N_d \\
    k' &= k \\
    k' + 1 &= k + N_d \\
    &\vdots
\end{align*}

Considering this shift, a transition of one step in the new time index $k'$ corresponds to a transition of $N_d$ 
steps in the original time index $k$. Therefore, the state at time $k'+1$ (which corresponds to $k+N_d$) can 
be expressed based on the state and control at time $k'$ (which is the same as $k$).

Let's consider the final prediction equation:
\noindent
\begin{equation}
    x\big[k+N_d\big] = A_b x\big[k\big] + B_b u\big[k\big]
\end{equation}
Using the new time index $k'$, this becomes:
\noindent
\begin{equation}
    x\big[k'+1\big] = A_b x\big[k'\big] + B_b u\big[k'\big]
\end{equation}
Now, defining the augmented state vector $X_a\big[k'\big]$ as:
\noindent
\begin{equation}
    X_a\big[k'\big] = 
    \begin{bmatrix}
        x\big[k'\big] \\
        z\big[k'-1\big]
    \end{bmatrix}
\end{equation}
with
\noindent
\begin{equation}
    z\big[k'-1\big] = u\big[k'\big]
\end{equation}
we get to the augmented state transition equation:
\noindent
\begin{equation}
    X_a \big[k'+1\big] = 
    \begin{bmatrix}
        A_b & B_b \\
        0 & 0
    \end{bmatrix}
    X_a\big[k'\big]
    +
    \begin{bmatrix}
        0 \\
        1
    \end{bmatrix}
    z\big[k'\big]
\end{equation}
where $z\big[k'\big] = u\big[k' + 1\big]$.

We can replace the augmented state transition matrix

\begin{equation}
    A_a = 
    \begin{bmatrix}
        A_b & B_b \\
        0 & 0
    \end{bmatrix}
\end{equation}

and the augmented input matrix

\begin{equation}
    B_a = 
    \begin{bmatrix}
        0 \\
        1
    \end{bmatrix}
\end{equation}

resulting in 

\begin{equation}
    X_a \big[k'+1\big] = 
    A_b X_a \big[k'\big] +
    B_a z\big[k'\big]
\end{equation}